% Generated draft: requires mathematical and source review.
\documentclass[11pt]{article}
\usepackage[margin=25mm]{geometry}
\usepackage{fontspec}
\setmainfont{Latin Modern Roman}
\usepackage{amsmath,amssymb,mathtools,unicode-math}
\setmathfont{Latin Modern Math}
\usepackage{enumitem,needspace,xurl,hyperref}
\setlength{\parindent}{0pt}
\setlength{\parskip}{5pt}
\setlength{\emergencystretch}{3em}
\newcommand{\N}{\mathbb{N}}
\newcommand{\Z}{\mathbb{Z}}
\newcommand{\Q}{\mathbb{Q}}
\newcommand{\R}{\mathbb{R}}
\newcommand{\C}{\mathbb{C}}
\newcommand{\F}{\mathbb{F}}
\newcommand{\A}{\mathbb{A}}
\newcommand{\PP}{\mathbb{P}}
\newcommand{\E}{\mathbb{E}}
\newcommand{\eps}{\varepsilon}
\newcommand{\dd}{\,\mathrm{d}}
\begin{document}
\section*{UnsolvedMath ID: 10000077}
\textbf{Source problem code:} catalogue entry
\subsection{Definitions and mathematical statement}
Let \((G_n,o_n)\) be finite planar graphs rooted at a uniform vertex and converging in distribution in rooted neighborhoods to \((G,o)\). Work with its infinite locally finite connected root component, when that is the intended hypothesis. For independent site percolation, with each vertex open with probability \(p\), define the quenched critical value
\[
p_c^{\mathrm{site}}(G)=\inf\{p:\mathbb P_p^G(o\leftrightarrow\infty)>0\}.
\]
The first requested statement is
\[
\Pr(p_c^{\mathrm{site}}(G)\ge1/2)=1.
\]
A literal graph-critical reading of the second is
\[
\text{for almost every }G,\qquad
\mathbb P_{p_c^{\mathrm{site}}(G)}^G
(\text{there exists an infinite open cluster})=0.
\]
This distinguishes randomness of the graph from percolation conditional on it.

A random rooted graph is unimodular if every nonnegative isomorphism-invariant measurable mass transport satisfies
\[
\mathbb E\sum_v F(G,o,v)=\mathbb E\sum_v F(G,v,o).
\]
The final question asks whether the same no-critical-cluster conclusion extends to every intended unimodular random graph, without silently imposing transitivity.

As written, the no-critical claim needs a qualification. Cycles converge locally to \(\Z\), a planar unimodular example with site threshold \(1\); at \(p=1\) its entire infinite graph is open. Thus the literal statement fails at a threshold-one case. A restriction \(p_c<1\), or a statement about the universal comparison parameter \(p=1/2\) rather than the graph's own threshold, would be a different version requiring source confirmation. The latter is compatible with the separate threshold lower bound but is not silently substituted. Infinity, local finiteness, degree/integrability assumptions, and quenched versus averaged formulations also require the intended scope.

The supplied research summary reports planar and unimodular partial progress but gives no precise theorem. Its broad claims of known cases for \(\Z^d\) and trees are not adopted without dimension/model qualifications; the threshold-one example already shows why those qualifiers matter.
\subsection{Short English statement}
Prove the planar-local-limit site threshold bound and clarify the intended absence of a critical infinite cluster, then ask for a unimodular extension. The literal no-critical statement includes the line at threshold one and therefore needs an explicit restriction or different parameter meaning.
\subsection{Sources}
\begin{itemize}
\item Local UnsolvedMath registry, permanent ID \(10000077\), exact source code \texttt{catalogue entry}, ``Critical percolation on distributional planar limits.'' Proposer Itai Benjamini.
\item Benjamini, \emph{Random Planar Metrics}, section 3 item (4), recorded PDF page 3; \url{https://www.wisdom.weizmann.ac.il/~itai/randomplanar2.pdf#page=3}. Recorded access 2026-07-29; extracted text checked against rendered source.
\item Recorded preservation copy: \url{https://arquivo.pt/noFrame/replay/20201231041544id_/http://www.wisdom.weizmann.ac.il/~itai/randomplanar2.pdf}.
\item Supplied PARTIAL-PROGRESS summary invokes Benjamini--Schramm planar-limit results and general critical-percolation work without complete references or precise assumptions.
\end{itemize}
\end{document}
